\documentclass[11pt]{article}
\usepackage[letterpaper,margin=1in]{geometry}
\usepackage{amsmath,amssymb,amsthm,mathtools}
\usepackage{booktabs,tabularx,array,microtype,xurl}
\usepackage[colorlinks=true,linkcolor=blue,citecolor=blue,urlcolor=blue]{hyperref}
\newtheorem{theorem}{Theorem}[section]
\newtheorem{proposition}[theorem]{Proposition}
\newtheorem{lemma}[theorem]{Lemma}
\newtheorem{corollary}[theorem]{Corollary}
\theoremstyle{remark}
\newtheorem{remark}[theorem]{Remark}
\newcommand{\R}{\mathbb R}
\newcommand{\C}{\mathbb C}
\newcommand{\E}{\mathbb E}
\newcommand{\dd}{\,\mathrm d}
\newcommand{\abs}[1]{\left\lvert#1\right\rvert}
\numberwithin{equation}{section}
\title{An explicit local zero-count bound\\for the Riemann heat deformation}
\author{Drafted for Edward Baker}
\date{17 October 2026\\\small Revised working draft; round-two extension awaiting cross-model review}
\begin{document}
\maketitle
\begin{abstract}
Polymath established local zero counts of order $\log X$ for the Riemann
heat deformation, uniformly for positive heat times at most $1/2$.
We extract explicit constants in the physical coordinate
$H_0(z)=\xi((1+iz)/2)/8$. The number of zeros with
$X<\operatorname{Re}z\le X+1$, counted with multiplicity, is at most
$(9/10)\log(X+2)+6$ for $0<t\le1/2$ and $X\ge10^5$.
For $X\ge2\cdot10^6$, the additive constant improves to $11/2$.
The proof makes an off-axis comparison uniform in time, extends a
Riemann--Siegel remainder estimate to both required real-part ranges,
and combines Schwarz--Pick with a conjugate-pair Poisson-kernel bound.
A separate height-selected remainder is retained in the truncation
estimate. All numerical constants and cutoffs are evaluated. The
lower-height range, cumulative error constants and an application gain
are not supplied. The new extension has same-model development checks
only; no specialist human review has occurred. No claim about RH or an
improved de Bruijn--Newman bound is made.
\end{abstract}

\noindent\textbf{Draft status.}
This revision integrates the requested round-two continuation of
17 October 2026. The earlier manuscript, at repository commit
\texttt{146b3f9}, received one cross-model LLM review (Claude Review 02,
recorded in commit \texttt{001aaeb}), accepting its results as proved from
named imports. That review does not cover the changed cutoff, sharper
extraction or this integrated revision. Those changes have same-model
checks only. The derivations, arithmetic and manuscript preparation used
\texttt{gpt-6.1-sol} (Codex), reasoning effort \texttt{ultra}, with
substantial LLM assistance. Compilation and layout checks are not
mathematical review. No literature-priority claim is made. The author
field identifies the intended recipient rather than asserting endorsement.

\section{Main result and relation to existing estimates}\label{sec:main}

Let
\begin{align}
 \Phi(u)&=\sum_{n\ge1}
 (2\pi^2n^4e^{9u}-3\pi n^2e^{5u})e^{-\pi n^2e^{4u}},\label{eq:phi}\\
 H_t(z)&=\int_0^\infty e^{tu^2}\Phi(u)\cos(zu)\dd u.
 \label{eq:heat}
\end{align}
The normalization is
\begin{equation}\label{eq:xi}
 H_0(z)=\tfrac18\xi((1+iz)/2),\qquad
 \xi(s)=\tfrac12s(s-1)\pi^{-s/2}\Gamma(s/2)\zeta(s),\qquad
 \partial_tH_t=-\partial_z^2H_t.
\end{equation}
Thus a zeta ordinate of magnitude $T$ corresponds to physical coordinate
$x=2T$. Every logarithm without a capital $\operatorname{Log}$ is natural.
For $X\ge0$, define
\begin{equation}\label{eq:counts}
 N_t(X)=\sum_{0<\operatorname{Re}\rho\le X}\operatorname{mult}(\rho),
 \qquad K_t(X)=N_t(X+1)-N_t(X),
\end{equation}
where the sum is over zeros of $H_t$. A zero on the right wall of a
window is included, one on its left wall is excluded, and a real double
zero contributes two.

\begin{theorem}[Explicit unit-window bound]\label{thm:count}
For $0<t\le1/2$ and $X\ge10^5$,
\begin{equation}\label{eq:main-count}
 K_t(X)\le\frac{2003}{2240}\log(X+2)+\frac{125}{21}
          \le\frac9{10}\log(X+2)+6.
\end{equation}
The count includes complex and multiple zeros; no reality, simplicity
or separation hypothesis is imposed.
\end{theorem}

\begin{corollary}[The original high-height domain]\label{cor:high-count}
For $0<t\le1/2$ and $X\ge2\cdot10^6$,
\begin{equation}\label{eq:high-count}
 K_t(X)\le\frac{2003}{2240}\log(X+2)+\frac{75}{14}
          \le\frac9{10}\log(X+2)+\frac{11}{2}.
\end{equation}
\end{corollary}

Theorem~1.5(iv) of Polymath~\cite{polymath} already establishes a
time-uniform estimate $K_t(X)=O(\log(2+X))$, together with a cumulative
estimate. The purpose here is to evaluate constants for one local
version on a declared domain. The cutoff is conservative and is not
claimed optimal. The range $200\le X<10^5$ remains outside the bound.
No consumer has supplied a numerical tolerance against which to measure
an application gain.

At $X=2\cdot10^6$, the previous exact allowance
$(82123/72000)\log(X+2)+574/27$ was about $37.81$ (its rounded version
$(6/5)\log(X+2)+22$ was $39.41$). The new rounded bounds are $19.06$
and $18.56$, respectively. These are upper allowances, not observed
counts. The leading mean density
$\log(X/(4\pi))/(4\pi)$ is about $0.953$ there, so the result is still
loose. Pairing conjugate zeros reduces the Poisson multiplier from
$41/9$ to $25/7$. Together with Schwarz--Pick, this replaces the earlier
additive allowance $574/27$ by $125/21$. These changes do not provide exact
counts or a lower population.

For a time-zero benchmark, Bellotti--Wong~\cite[Corollary 1.6]{bellotti-wong}
give
\[
 N_\zeta(T+1)-N_\zeta(T)\le\frac{\log T}{2\pi}+4.8405
 \qquad(3\le T+1\le30\,610\,046\,000).
\]
At $T=10^6$ the right side is less than $7.04$. A unit physical window
corresponds to a half-unit $T$ window and is contained in this unit
$T$ window. This is context at $t=0$, not an input to the proof or an
estimate transferred to positive heat time.

The analytic input is a two-regime comparison. The nonvanishing carrier
$M_t$ is defined in Section~\ref{sec:inputs}.
\begin{theorem}[Uniform off-axis comparison]\label{thm:comparison}
For $0<t\le1/2$, $x\ge10^5$ and $3\le y\le4$,
\begin{equation}\label{eq:comparison}
 \abs{\frac{H_t(x+iy)}{M_t((1+y-ix)/2)}-1}
 <\frac{5094127}{7350000}<\frac{347}{500}<\frac7{10}.
\end{equation}
If in addition $x\ge2\cdot10^6$, then
\begin{equation}\label{eq:comparison-high}
 \abs{\frac{H_t(x+iy)}{M_t((1+y-ix)/2)}-1}
 <q_0:=\frac{96624791}{147000000}<\frac{48}{73}<\frac{33}{50}.
\end{equation}
\end{theorem}

Section~\ref{sec:budget} proves both comparisons.
Section~\ref{sec:count} then obtains the counts using Schwarz--Pick
and a positive logarithmic-derivative sum with conjugate pairs.
No compactness constant or finite zero survey enters the argument.

\section{Carriers and imported analytic estimates}\label{sec:inputs}

On $\C\setminus(-\infty,1]$, with principal logarithms, set
\begin{align}
 M_0(s)&=\frac{s(s-1)}{16}\pi^{-s/2}\sqrt{2\pi}
   \exp\left\{\left(\frac s2-\frac12\right)
                        \operatorname{Log}\frac s2-\frac s2\right\},
   \label{eq:M0}\\
 \alpha(s)&=\frac{M_0'(s)}{M_0(s)}
   =\frac1{2s}+\frac1{s-1}+\frac12\operatorname{Log}\frac{s}{2\pi},
   \label{eq:alpha}\\
 M_t(s)&=e^{t\alpha(s)^2/4}M_0(s),\qquad
 \frac{M_t'}{M_t}(s)=\alpha(s)+\frac t2\alpha(s)\alpha'(s).
   \label{eq:Mt}
\end{align}
The functions have the expected conjugation symmetry and are nonzero
on the domains used below. We use the conventions
\begin{equation}\label{eq:parameters}
 T=\frac x2,\quad T'=T+\frac{\pi t}{8},\quad
 Q=\frac{x}{4\pi},\quad L=\log Q,\quad
 a=\sqrt{\frac{T'}{2\pi}},\quad N=\lfloor a\rfloor,
\end{equation}
and
\begin{equation}\label{eq:branches}
 s_+=\frac{1+y-ix}{2},\qquad s_-=\frac{1-y+ix}{2},\qquad
 b_n^t=e^{t\log^2 n/4},\qquad B_t(z)=M_t((1-iz)/2).
\end{equation}
For $3\le y\le4$, the real parts of $s_\pm$ belong to
$[-3/2,5/2]$. Note that $a^2=Q+t/16$.

We state the imported ingredients precisely. All equation and theorem
numbers attributed to Polymath refer to arXiv version~2 of
\cite{polymath}. First, $H_t$ is even and real entire of order one,
$H_t(iy)>0$ for real $y$, and its zeros lie in $|\operatorname{Im}z|\le1$
for $0\le t\le1/2$; see its Section~3 and Theorem~3.2.

To describe its expansion~(39), write
\[
 c(s)=\frac{s(s-1)}{16}\pi^{-s/2}\Gamma(s/2),\qquad
 r_{0,n}(s)=c(s)n^{-s},
\]
and
\[
 R_{0,N}(s)=c(s)\int_{N\swarrow N+1}
       \frac{w^{-s}e^{i\pi w^2}}{e^{\pi iw}-e^{-\pi iw}}\dd w.
\]
The contour has direction $e^{5\pi i/4}$ and crosses $(N,N+1)$.
For nonreal $s$, define $r_{t,n}$ and $R_{t,N}$ by Gaussian convolution
in $s$, for example
\begin{equation}\label{eq:Rheat}
 R_{t,N}(s)=\frac1{\sqrt\pi}\int_\R
                  R_{0,N}(s+\sqrt t\,v)e^{-v^2}\dd v.
\end{equation}
For $F^*(s)=\overline{F(\overline s)}$, the exact formula is
\begin{equation}\label{eq:exact-expansion}
 \begin{split}
 H_t(z)={}&\sum_{n=1}^N r_{t,n}(s_-)
          +\sum_{n=1}^N r_{t,n}^*(s_+)\\
          &+R_{t,N}(s_-)+R_{t,N}^*(s_+).
 \end{split}
\end{equation}
This formula and its contour shifts are imported from
\cite[(39)--(41)]{polymath}; sub-Gaussian growth on the shifted lines
justifies the Gaussian integrals.

Proposition~6.1 of \cite{polymath} applies for every real $\sigma$,
$T>10$, $n\ge1$ and $0<t\le1/2$. It gives
\begin{align}
 r_{t,n}(s)&=M_t(s)b_n^t n^{-s-(t/2)\alpha(s)}(1+\theta_{t,n}(s)),
 \quad s=\sigma+iT,\label{eq:r-approx}\\
 \abs{\theta_{t,n}(s)}&\le
 \exp\left(\frac{(t^2/8)|\alpha(s)-\log n|^2+t/4+1/6}{T-3.33}\right)-1.
 \label{eq:r-error}
\end{align}
The $r$ functions have conjugation symmetry, so the same estimates hold
at the reflected arguments. We also import the general-real-part
coefficient and remainder bounds of Proposition~6.2, including
$|C_0(p)|\le1/2$, and the explicit Stirling estimate of Lemma~5.1(v).
We do not extend the final constants of its Proposition~6.3 beyond
their stated real-part range. Section~\ref{sec:remainder} rebuilds the
needed remainder estimate instead. The coefficient and remainder bounds,
including the $C_0$ estimate, originate with Arias de Reyna, and the
Stirling estimate with Boyd, as cited by Polymath. We import those bounds
through Polymath; their original papers have not been checked here.

\section{Bounds for the ordinary terms}\label{sec:ordinary}

\begin{lemma}[Heat weights and termwise error]\label{lem:ordinary}
For $x\ge10^5$, $0<t\le1/2$, $-3/2\le\sigma\le5/2$ and
$s=\sigma\pm iT$, uniformly for $1\le n\le N$,
\begin{align}
 \operatorname{Re}\alpha(s)&\ge L/2-13/x^2,\label{eq:alpha-real}\\
 \abs{\alpha(s)-L/2\mp i\pi/4}
    &\le\frac{11}{2x}+\frac{25}{4x^2}<\frac6x,\label{eq:alpha-small}\\
 W_n(s):=b_n^t n^{-(t/2)\operatorname{Re}\alpha(s)}&\le1,
 \qquad |\theta_{t,n}(s)|<\frac1{6000}.\label{eq:weights-error}
\end{align}
For $x\ge2\cdot10^6$, the termwise error improves to $1/50000$.
\end{lemma}
\begin{proof}
In \eqref{eq:alpha}, the real rational terms are bounded below by
$-13/x^2$, and the log-modulus correction to $L/2$ is nonnegative.
The inequalities $|\arctan(\sigma/T)|\le|\sigma|/T$ and
$\log(1+\sigma^2/T^2)\le\sigma^2/T^2$, together with the rational
terms, prove \eqref{eq:alpha-small}. Since $a^2=Q+t/16$,
\[
 \log N\le L/2+\frac{\pi}{16x},\qquad
 \abs{\alpha(s)-\log n}\le L/2+1\le\log x.
\]
The exact identity
\begin{equation}\label{eq:weight-identity}
 W_n(s)=\exp\left\{\frac t4\log n
                  (\log n-2\operatorname{Re}\alpha(s))\right\}
\end{equation}
and the bound
$2\operatorname{Re}\alpha(s)-\log N
 \ge L/2-\pi/(16x)-26/x^2>0$ prove $W_n\le1$.
The heat coefficient in \eqref{eq:weight-identity} is $t/4$.

Because $T-3.33\ge x/3$, the exponent in \eqref{eq:r-error} is at most
\[
 u(x)=\frac{3(\log x)^2/32+7/8}{x}.
\]
Both $(\log x)^2/x$ and $1/x$ decrease here. Since $\log10^5<12$,
\[
 u(x)\le\frac{23}{160000}<\frac18,\qquad
 e^{u(x)}-1\le\frac87u(x)\le\frac{23}{140000}<\frac1{6000}.
\]
For $x\ge2\cdot10^6$, use $\log(2\cdot10^6)<15$ to obtain
$u(x)\le703/(32\cdot2\cdot10^6)$ and
$e^{u(x)}-1<1/50000$. This proves both assertions.
\end{proof}

The dominant branch has real part $(1+y)/2\ge2$. Its leading
normalized approximant is exactly $1$, and its other terms total at most
\begin{equation}\label{eq:main-excess}
 \sum_{n=2}^N n^{-(1+y)/2}W_n(s_+)
 \le\zeta(2)-1\le\frac{95}{147}.
\end{equation}
Here and below $\zeta(2)\le(22/7)^2/6=242/147$ is used as a convenient
rational upper bound.

\section{A remainder estimate for both real-part ranges}\label{sec:remainder}


\begin{lemma}[A remainder bound on an enlarged real-part interval]
\label{lem:remainder}
Let
\[
 T\ge 50000,\qquad -\frac32\le\sigma\le\frac52,
 \qquad 0<t\le\frac12,
\]
and set
\[
 T'=T+\frac{\pi t}{8},\qquad
 a=\sqrt{\frac{T'}{2\pi}},\qquad N=\lfloor a\rfloor.
\]
For the remainder $R_{t,N}$ defined in \cite[equations (38)--(41)]{polymath},
one has
\begin{equation}
 \label{eq:remainder-expanded-domain}
 \bigl|R_{t,N}(\sigma+iT)\bigr|
 <2\exp\!\left(\frac{t\pi^2}{64}\right)|M_0(iT')|.
\end{equation}
For $T\ge10^6$, the prefactor $2$ can be replaced by $1$:
\begin{equation}\label{eq:remainder-high}
 |R_{t,N}(\sigma+iT)|\le e^{t\pi^2/64}|M_0(iT')|.
\end{equation}
\end{lemma}

\begin{proof}
We use the general-real-part estimates of
\cite[Proposition 6.2]{polymath}, rather than extending the stated domain
of its Proposition 6.3.  Write
\[
 D=T-3,\qquad p=1-2(a-N),\qquad
 u=\sigma+\sqrt t\,v.
\]
Expectation below refers to the density
\[
 \frac{1}{\sqrt{\pi t}}
 \exp\!\left(-\frac{(u-\sigma)^2}{t}\right)\,du.
\]
Throughout the proof, $D\ge49997$ and
$a^2>87500/11>89^2$. The larger endpoint $a>300$ is used only for the
final high-height subcase.

\emph{The Gaussian envelope.}
Use the shift $\beta=\pi i/4$ in
\cite[equation (41)]{polymath}.  Its imaginary displacement is
$\pi t/8$, so both the original and shifted contours remain in the
upper half-plane.  Apply Proposition 6.2 to $u+iT'$.
After extracting the common carrier and its unit-modulus phase, the
integrand is
\[
 e^{\eta(u)}\left(
 C_0(p)+\sum_{k=1}^{K_u}\frac{C_k(p,u)}{a^k}
                    +RS_{K_u}(u+iT')\right),
\]
where $K_u$ will be chosen below.  The extracted carrier has modulus
\[
 \exp\!\left(\frac{t\pi^2}{64}\right)|M_0(iT')|.
\]
The explicit Stirling estimate in \cite[Lemma 5.1(v)]{polymath}, the
formula for $\alpha$ in its equation (9), and the derivative bound in
its equation (43) give
\begin{align}
 |\eta(u)|
 &\le \frac{u^2}{4(T'-3)}
       +\frac{3|u|}{2T'}+\frac{1}{6(T'-0.66)}\notag\\
 &\le\frac{u^2+6|u|+2/3}{4D}
 \le\frac{u^2+1}{D}.
 \label{eq:remainder-gaussian-envelope}
\end{align}
Indeed, the horizontal Taylor segment from $iT'$ to $u+iT'$ has
constant imaginary part $T'$, and
$6|u|\le3u^2+3$.  Thus this estimate is valid for every real $u$;
no boundedness assumption on $u$ enters the Taylor argument.

For a real number $b$, put
\[
 J_b=\mathbb E\exp\!\left(\frac{u^2+1}{D}+bu\right),
 \qquad h=1-\frac tD.
\]
Completing the square gives the exact formula
\begin{equation}
 \label{eq:remainder-gaussian-moment}
 J_b=h^{-1/2}
 \exp\!\left(\frac{\sigma^2+1}{D}+b\sigma
       +\frac{t(\sigma/D+b/2)^2}{h}\right).
\end{equation}
The simpler inequality $u^2\le2\sigma^2+2tv^2$ also gives
\begin{equation}
 \label{eq:remainder-zero-moment}
 J_0\le
 \exp\!\left(\frac{2\sigma^2+t+1}{T-6}\right)
 \le\exp\!\left(\frac{14}{T-6}\right)<\frac43.
\end{equation}
For $b=\log9$ or $b=-\log2$, the factor in
\eqref{eq:remainder-gaussian-moment} remaining after removal of
$e^{b\sigma}$ is less than $2$.  To see this, use
\[
 h\ge\frac{999}{1000},\qquad
 \frac{\sigma^2+1}{D}<\frac1{1000},\qquad
 \frac{|\sigma|}{D}<\frac1{1000},\qquad
 -\frac12\log h<\frac1{1000}.
\]
Since $\log3<11/10$ and $\log\sqrt2<\log3$, its logarithm is less than
\[
 \frac2{1000}
 +\frac12\left(\frac{1101}{1000}\right)^2\frac{1000}{999}
 <\frac23<\log2.
\]

\emph{Coefficient summation and the fixed-index remainder.}
Choose
\[
 m=\lfloor2a^2\rfloor,\qquad
 K_u=\begin{cases}
 1,&u\ge0,\\
 \max(\lfloor-u\rfloor+3,m),&u<0.
 \end{cases}
\]
When $u<0$, this choice satisfies $K_u+u\ge2$, as required by
\cite[equation (57)]{polymath}.  Define
\[
 S(u)=\sum_{k=1}^{K_u}\frac{|C_k(p,u)|}{a^k}
                     +|RS_{K_u}(u+iT')|,
 \qquad E_a=a^{-1}e^{-a^2/2}.
\]
The following pointwise bound includes every remainder term:
\begin{align}
 S(u)\le{}&
 \frac{0.2\,9^u}{a-0.865}
 +\frac{0.029\,2^{-u}}{a-0.353}\notag\\
 &+\frac{2^{-u}}2
       \sum_{2a^2\le k\le-u+4}
          \frac{1.1^k\Gamma(k/2)}{a^k}
 +E_a.
 \label{eq:remainder-pointwise-majorant}
\end{align}
All sums over $k$ have integer indices; an empty sum is zero.

Here are the details of the summation.  For $u\ge0$, the bounds in
\cite[equations (54) and (56)]{polymath} give
\[
 S(u)\le\frac{0.200\,9^u}{a}
          +\frac{0.173\,2^{3u/2}}{a^2}
 \le\frac{0.2\,9^u}{a-0.865}.
\]
The value $u=0$ can be assigned by continuity and has no effect on
the integrals.  For $u<0$, the coefficients with $k\le m$ satisfy
\[
 \sum_{k=1}^{m}\frac{|C_k(p,u)|}{a^k}
 \le 0.036\,2^{-u}
 \left(\frac{0.445\sqrt\pi}{a}
       +\frac{0.247}{a^2}+\frac{0.098}{a^3}\right)
 \le\frac{0.029\,2^{-u}}{a-0.353}.
\]
This follows from \cite[equation (55)]{polymath} by separating even
and odd indices and using
\[
 \frac{\Gamma((k+2)/2)}{a^{k+2}}
 =\frac{k}{2a^2}\frac{\Gamma(k/2)}{a^k}
 \le\frac{\Gamma(k/2)}{a^k}
 \qquad(k\le2a^2).
\]
The constants bound the corresponding geometric sums. The final
compression remains valid at the smaller endpoint, since
\[
 .036\left(.445\frac{71}{40}+
              \frac{.247}{89}+\frac{.098}{89^2}\right)<.029,
 \qquad \sqrt\pi<71/40.
\]

If $K_u>m$, every coefficient beyond $m$ has index greater than
$2a^2$.  Moreover the remainder index satisfies
\[
 K_u+1=\lfloor-u\rfloor+4>2a^2,
 \qquad K_u+1\le-u+4.
\]
The bounds in equations (55) and (57) of \cite{polymath} therefore
place these additional coefficients and the remainder inside the
tail sum in \eqref{eq:remainder-pointwise-majorant}.

If $K_u=m$, the remainder need not belong to that tail.  To bound it,
put $\nu=(m+1)/2$, so that $a^2<\nu\le a^2+1/2$.
Equation (57) and explicit real Stirling give
\begin{align*}
 |RS_m(u+iT')|
 &\le\frac12\Gamma(\nu)(1.1/a)^{2\nu}\\
 &\le\frac3{2a}
     \exp\!\left(\nu\log(\nu/a^2)-\frac45\nu\right)\\
 &<\frac5a e^{-4a^2/5}\le E_a.
\end{align*}
We used
$\Gamma(\nu)\le3\nu^{\nu-1/2}e^{-\nu}$,
$\log(1.21)<1/5$, $\log(\nu/a^2)\le1/(2a^2)$,
$\nu\log(\nu/a^2)\le1/2+1/(4a^2)<1$, and
$3a^2/10>2>\log5$. These bounds hold already for $a>89$.
Thus the added $E_a$ accounts for the fixed-index
remainder even when the displayed tail range is empty.  This proves
\eqref{eq:remainder-pointwise-majorant} in every case.

\emph{Integration of the four contributions.}
Let $\delta_1,\ldots,\delta_4$ denote the expectations of the four
terms in \eqref{eq:remainder-pointwise-majorant}, each multiplied by
$\exp((u^2+1)/D)$.  The moment bounds above give
\begin{align}
 \delta_1+\delta_2
 &<\frac{0.4\,9^\sigma+0.058\,2^{-\sigma}}{a-0.865}\notag\\
 &\le\frac{97.374}{a-0.865}.
 \label{eq:remainder-first-two-deltas}
\end{align}
Here $9^\sigma\le243$ and $2^{-\sigma}<3$. Retain this denominator:
the shortcut $97.374/(a-.865)\le98/a$ requires $a\ge42385/313$
and is unavailable at $a=89$.

In the third contribution, $2a^2>15842$, so enlarge the integer
summation to $k\ge14$.
Its integration region has $u\le4-k\le-10$.  Since $|\sigma|\le5/2$,
\[
 (u-\sigma)^2\ge\frac9{16}u^2,
 \qquad
 \frac{u^2+1}{D}-\frac{(u-\sigma)^2}{t}
 \le\frac1D-\frac{u^2}{2t}.
\]
Furthermore,
\[
 \int_{-\infty}^{4-k}e^{-u^2/(2t)-u\log2}\,du
 \le\frac{2t}{k-6}
       e^{-(k-4)^2+(k-4)\log2}.
\]
For
\[
 f(k)=1.1^k\Gamma(k/2)
             e^{-(k-4)^2+(k-4)\log2},
\]
one has $f(14),f(15)<10^{-30}$ and
\[
 \frac{f(k+2)}{f(k)}=2.42k\,e^{-4k+12}<1
 \qquad(k\ge14).
\]
The initial bounds follow, for example, from
$\Gamma(7)=720$, $\Gamma(15/2)=135135\sqrt\pi/128$,
$\sqrt\pi<2$, and $e>8/3$.  Summing the resulting geometric series
in $1/a$ gives
\begin{equation}
 \label{eq:remainder-tail-delta}
 \delta_3\le\frac{4\cdot10^{-30}}{a^{14}}<\frac1a.
\end{equation}
The fourth contribution satisfies
\begin{equation}
 \label{eq:remainder-fixed-index-delta}
 \delta_4=E_aJ_0<\frac43 E_a<\frac1a.
\end{equation}

Finally, \cite[equation (58)]{polymath} gives $|C_0(p)|\le1/2$.
Taking absolute values in the shifted integral and using
\eqref{eq:remainder-gaussian-envelope} therefore yields
\[
 \frac{|R_{t,N}(\sigma+iT)|}
 {e^{t\pi^2/64}|M_0(iT')|}
 \le\frac12J_0+\sum_{j=1}^{4}\delta_j
 <\frac23+\frac{97.374}{a-.865}+\frac2a.
\]
The last expression decreases in $a$ and is less than
\[
 \frac23+\frac{97.374}{89-.865}+\frac2{89}
 =\frac{42215698}{23532045}<\frac95<2.
\]
This proves \eqref{eq:remainder-expanded-domain}.
For $T\ge10^6$, one has $a>300$, and now
$97.374/(a-.865)\le98/a$ is permitted. The same strict chain yields
\[
 \frac{|R_{t,N}(\sigma+iT)|}{e^{t\pi^2/64}|M_0(iT')|}
 <\frac12\cdot\frac43+\frac{100}{300}=1,
\]
proving \eqref{eq:remainder-high}. The endpoint on the right is an
equality of rational allowances; strictness comes from the preceding
analytic inequalities.
\end{proof}

\section{Ratios of the carriers}\label{sec:carriers}

\begin{lemma}[Carrier comparisons]\label{lem:carriers}
On the domain of Theorem~\ref{thm:comparison},
\begin{align}
 \abs{\frac{M_t(s_-)}{M_t(s_+)}}&<2Q^{-y/2},\label{eq:gamma}\\
 e^{t\pi^2/64}\frac{|M_0(iT')|}{|M_t(s_+)|}
   &\le Q^{-(1+y)/4}.\label{eq:remainder-carrier}
\end{align}
\end{lemma}
\begin{proof}
Put $A=(1+y)/2\in[2,5/2]$ and $C=1-A\in[-3/2,-1]$.
For $-3/2\le\sigma\le5/2$, direct use of \eqref{eq:alpha} gives
\begin{align*}
 \abs{\operatorname{Re}\alpha(\sigma+iT)-L/2}&\le85/(16T^2),&
 \abs{\operatorname{Im}\alpha(\sigma+iT)-\pi/4}&\le11/(4T),\\
 |\alpha(\sigma+iT)|&\le L/2+1,&
 |\alpha'(\sigma+iT)|&\le1/(2T)+3/(2T^2).
\end{align*}
Integrating the real part of the logarithmic derivative in
\eqref{eq:Mt} from $C+iT$ to $A+iT$ gives
\[
 \log\abs{\frac{M_t(s_-)}{M_t(s_+)}}
 \le-\frac y2L+\frac{85}{x^2}
                 +(L/2+1)(1/x+6/x^2)
 \le-\frac y2L+\frac{\log x}{x}.
\]
For example, the final inequality follows from $L\le\log x$,
$85/x\le1/100$, $6/x\le1/100$ and $\log x\ge3$.
Conjugation equates the denominator's modulus to that at $A+iT$.
Since $\log x/x<1/2$ and $e^{\log x/x}<2$, this proves
\eqref{eq:gamma}.

For the second inequality, set
$G(\sigma,\tau)=\operatorname{Re}\log M_0(\sigma+i\tau)$.
An exact horizontal integration gives
\begin{align*}
 G(A,T)-G(0,T)-AL/2
   ={}&\tfrac14\log(1+A^2/T^2)
       +\tfrac12\log\frac{T^2+(A-1)^2}{T^2+1}\\
      &+\tfrac14\int_0^A\log(1+\sigma^2/T^2)\dd\sigma\ge0.
\end{align*}
Every term is nonnegative because $A\ge2$. On the imaginary axis,
\[
 \partial_\tau G(0,\tau)
 =-\frac\pi4+\frac1{2\tau}+\frac{\tau}{1+\tau^2}
 \le-\frac\pi4+\frac3{2T}\qquad(\tau\ge T).
\]
At $A+iT$, the exact formula for $\alpha$ gives
$\operatorname{Re}\alpha\ge L/2$ and
$0\le\operatorname{Im}\alpha\le\pi/4$. Thus
\[
 \operatorname{Re}(\alpha(A+iT)^2)\ge L^2/4-\pi^2/16.
\]
Combining these estimates, including the displacement
$T'-T=\pi t/8$, yields
\begin{equation}\label{eq:carrier-cancellation}
 e^{t\pi^2/64}\frac{|M_0(iT')|}{|M_t(s_+)|}
 \le Q^{-A/2}\exp\left(-\frac{tL^2}{16}+\frac{3\pi t}{8x}\right).
\end{equation}
The exponent is nonpositive because $L>1>6\pi/x$ on the stated
domain. This proves \eqref{eq:remainder-carrier}.
\end{proof}

\section{The complete off-axis budget}\label{sec:budget}

\begin{proof}[Proof of Theorem~\ref{thm:comparison}]
Use $N=\lfloor\sqrt{Q+t/16}\rfloor$ in the exact expansion.
Lemma~\ref{lem:ordinary} bounds the dominant excess by $95/147$.
For the opposite branch,
\[
 N^{(y+1)/2}\le(Q+t/16)^{(y+1)/4}<2Q^{(y+1)/4},
\]
because $(1+t/(16Q))^{(y+1)/4}<(33/32)^2<2$. Hence
\begin{equation}\label{eq:opposite}
 \abs{\frac{M_t(s_-)}{M_t(s_+)}}
 \sum_{n=1}^N n^{(y-1)/2}W_n(s_-)
 <4Q^{(1-y)/4}\le4Q^{-1/2}<\frac{23}{500}.
\end{equation}
The last inequality follows from $4\pi<13$ and
$16\cdot13/10^5<(23/500)^2$.
The normalized ordinary approximation errors are less than
\[
 \frac{242/147+23/500}{6000}<\frac3{10000}.
\]
Lemma~\ref{lem:remainder} applies because $T\ge50000$. Its prefactor
$2$ and the carrier ratio bound both remainders by
\[
 4Q^{-(1+y)/4}\le4/Q<52/x.
\]
Thus
\begin{align*}
 \abs{H_t(x+iy)/M_t(s_+)-1}
 &<\frac{95}{147}+\frac{23}{500}+\frac3{10000}+\frac{52}{100000}\\
 &=\frac{5094127}{7350000}<\frac{347}{500}<\frac7{10}.
\end{align*}

For $x\ge2\cdot10^6$, retain the sharper ordinary error $1/50000$
and the remainder prefactor $1$ in \eqref{eq:remainder-high}.
The opposite branch is then less than $11/1000$, since
$16\cdot13/(2\cdot10^6)<(11/1000)^2$.
The ordinary errors total less than $1/25000$, and the two remainders
total at most $2/Q<26/x$. Therefore
\begin{align*}
 \abs{H_t(x+iy)/M_t(s_+)-1}
 &<\frac{95}{147}+\frac{11}{1000}+\frac1{25000}+\frac{26}{2\cdot10^6}\\
 &=\frac{96624791}{147000000}<\frac{48}{73}<\frac{33}{50}.
\end{align*}
All estimates are uniform on their respective stated domains.
\end{proof}

\section{From the comparison to a local count}\label{sec:count}

\begin{lemma}[A disk logarithmic-derivative estimate]\label{lem:disk}
Let $F$ be holomorphic near a closed disk $|z-w|\le r$ and satisfy
$|F-1|\le q$, where $r>0$ and $0<q<1$. Then
\begin{equation}\label{eq:disk}
 \abs{\frac{F'(w)}{F(w)}}
 \le\frac{2q}{r(1+\sqrt{1-q^2})}.
\end{equation}
\end{lemma}
\begin{proof}
Set $f(\zeta)=(F(w+r\zeta)-1)/q$ and $a=f(0)$.
Schwarz--Pick gives $|f'(0)|\le1-|a|^2$; for a constant map with
$|a|=1$, the same statement holds with derivative zero. With $s=|a|$,
\[
 \abs{F'(w)/F(w)}
 \le\frac q r\frac{1-s^2}{1-qs}
 \le\frac{2q}{r(1+\sqrt{1-q^2})}.
\]
The last maximum occurs at $s=q/(1+\sqrt{1-q^2})$, as follows by
differentiation. All denominators are positive for $0\le s\le1$.
\end{proof}

\begin{lemma}[Comparison-to-count extraction]\label{lem:extract}
Fix $0<t\le1/2$ and $X\ge200$, and put
\[
 U=X+\tfrac12,\qquad w=U+\tfrac72i,\qquad F_t=H_t/B_t.
\]
If $\sup_{|z-w|\le1/2}|F_t(z)-1|\le q$ for $0<q<1$, then
\begin{equation}\label{eq:generic-count}
 K_t(X)\le\frac{2003}{2240}\log(X+2)
        +\frac{25}{7}\frac{4q}{1+\sqrt{1-q^2}}.
\end{equation}
\end{lemma}
\begin{proof}
The disk has imaginary parts in $[3,4]$ and maps under $s=(1-iz)/2$
to real part at least $2$. Thus $B_t$ is nonzero and $F_t$ is
holomorphic near the closed disk. Lemma~\ref{lem:disk} applies with
$r=1/2$. At $s=(1-iw)/2=9/4-iU/2$, let $L_U=\log(U/(4\pi))$.
The exact formula for $\alpha$ gives
\begin{align*}
 |\alpha(s)|&\le L_U/2+\pi/4+3/U+81/(16U^2)
              <L_U/2+1<\tfrac12\log(X+2),\\
 |\alpha'(s)|&\le1/U+6/U^2\le103/20000<3/500.
\end{align*}
Here $4\pi>e^2$ and $U<X+2$. Holding time fixed,
\begin{equation}\label{eq:Bprime}
 \frac{B_t'}{B_t}(w)=-\frac i2\alpha(s)
    \left(1+\frac t2\alpha'(s)\right),\qquad
 \abs{B_t'/B_t}(w)\le\frac{2003}{8000}\log(X+2).
\end{equation}
Evenness, order one and $H_t(0)>0$ give the paired Hadamard identity
\cite[(105)]{polymath}. Put $v=7/2$. Its imaginary part is
\begin{equation}\label{eq:poisson}
 D_t(U):=-\operatorname{Im}\frac{H_t'}{H_t}(w)
 =\sum_\rho\frac{v-\operatorname{Im}\rho}
 {(U-\operatorname{Re}\rho)^2+(v-\operatorname{Im}\rho)^2}.
\end{equation}
The scalar sum is positive and absolutely convergent by the strip
bound and order-one growth, independently of the count being proved.

For a zero in the unit window put $h=U-\operatorname{Re}\rho$;
then $|h|\le1/2$. Set $p_h(a)=a/(a^2+h^2)$. On $[5/2,9/2]$,
\[
 p_h''(a)=\frac{2a(a^2-3h^2)}{(a^2+h^2)^3}>0.
\]
The real source makes $H_t$ real entire, so conjugate nonreal zeros
have equal multiplicity and the same window membership. If their
imaginary parts are $\pm b$, the strip bound $|b|\le1$ and convexity give
\[
 \frac{p_h(v-b)+p_h(v+b)}2\ge p_h(v)
 \ge\frac{7/2}{(7/2)^2+1/4}=\frac7{25}.
\]
Real zeros obey the same per-multiplicity inequality directly. Thus
\begin{equation}\label{eq:paired-kernel-count}
 K_t(X)\le\frac{25}{7}D_t(U).
\end{equation}
Combining the disk estimate, \eqref{eq:Bprime} and
\eqref{eq:poisson} proves \eqref{eq:generic-count}.
Pairing preserves membership because the window restricts only real
parts. Both walls satisfy the distance bound, so the specified half-open
convention and multiple roots require no further exception.
\end{proof}

\begin{proof}[Proof of Theorem~\ref{thm:count}]
For $X\ge10^5$, the whole extraction disk lies in
Theorem~\ref{thm:comparison}'s domain. Use $q=7/10$.
The derivative allowance in \eqref{eq:disk} is less than $5/3$:
for $0\le s\le1$, clearing its positive denominator reduces
\[
 \frac75\frac{1-s^2}{1-(7/10)s}\le\frac53
\]
to $84s^2-70s+16=84(s-5/12)^2+17/12>0$.
Lemma~\ref{lem:extract} consequently gives
\[
 K_t(X)\le\frac{2003}{2240}\log(X+2)+\frac{125}{21}.
\]
Since $2003/2240<9/10$ and $125/21<6$, the rounded assertion follows.
\end{proof}

\begin{proof}[Proof of Corollary~\ref{cor:high-count}]
For $X\ge2\cdot10^6$, use $q_0<48/73$ from
\eqref{eq:comparison-high}. The disk allowance is increasing in $q$,
and at $q=48/73$ its square root is $55/73$ and
\[
 \frac{4q}{1+\sqrt{1-q^2}}=\frac32.
\]
The additive count allowance is therefore $75/14<11/2$.
\end{proof}

\section{Scope and limitations}\label{sec:scope}

This is an explicit upper count on a restricted high-height domain.
An upper count alone supplies neither a lower population, an exact
interior integer count, nor a lower bound on a derivative at a zero.
The theorem is compatible with multiple zeros and collisions. The
paired kernel argument applies to any real entire function meeting its
Hadamard and strip hypotheses; the complete theta expansion is the additional input used
to obtain the particular numerical comparison. Positivity, evenness or
reflection of a modified source does not supply that expansion.

In particular, the Euler-factor, weighted-square, reflected and
positive-primitive control families recorded in the research repository
do not challenge a multiplicity-count upper bound merely by colliding.
No untested control is claimed to fail \eqref{eq:comparison} on a
particular box. No noncollision conclusion is drawn.

The lower-height gap is stated rather than filled by compactness.
The local argument does not provide an evaluated cumulative estimate
for $N_t(X)$. If a later application uses Polymath's cumulative
comparison function, the normalization to retain is
\[
 g(X,t)=\frac{X}{4\pi}\log\frac{X}{4\pi}-\frac{X}{4\pi}
        +\frac{11}{8}+\frac{t}{16}\log\frac{X}{4\pi}.
\]
The heat drift and the constant $11/8$ should not be replaced by a
time-zero or averaged convention. A conductor-dependent version for
Dirichlet families also requires additional estimates.

The remaining height gap is $200\le X<10^5$. The bounded continuation
stops at its present result, pending cross-model review of the changed
cutoff and extraction. A separate application must specify its kernel,
domain, baseline and numerical tolerance before further optimization or
a gain assessment. In particular, a principal-value velocity tail needs
symmetric cancellation and two-sided cumulative information; an upper
local count alone does not provide these. The draft makes no assertion
that the bound improves an existing certificate.

\appendix
\section{Source normalization and reproducibility}\label{sec:replay}

The source statements actually used are identified in
Section~\ref{sec:inputs}. In the proof following (99) of
\cite{polymath}, the displayed logarithmic contribution to
$\operatorname{Re}s_*$ uses $t/2$, while its definitions (9), (17)
and its explicit estimate (21) give $t/4$. We avoid that display and
use the exact identity \eqref{eq:weight-identity}. This observation
does not dispute the source's time-uniform asymptotic theorem.

Immediately before (60) of \cite{polymath}, the printed inequality is
\[
 \frac1{6(T'-.66)}\le\frac{1/3}{4(T-3)}.
\]
It fails by a factor tending to two. Replacing the numerator $1/3$ by
$2/3$ gives the valid bound used in
\eqref{eq:remainder-gaussian-envelope}. Our envelope $(u^2+1)/D$ is
conservative. Also, the displayed tail for (63) can omit the remainder
when the height index wins: its upper index range may be empty although
$RS_m$ is still present. The separate $E_a$ in
\eqref{eq:remainder-pointwise-majorant} retains it. These small source
slips do not defeat Polymath's conclusions.

The standard-library arithmetic programs in \path{../numerics/} are:
\begin{itemize}
\item \path{check_local_count_reduction.py}: the original kernel,
      carrier chain rule and extraction constants.
\item \path{check_off_axis_comparison.py}: the original high-height
      comparison budget and Gaussian/remainder allowances.
\item \path{round-2/check_lower_cutoff_and_paired_extraction.py}:
      the reduced cutoff, retained full remainder denominator,
      Schwarz--Pick constants and conjugate-pair kernel.
\end{itemize}
Their small JSON records identify the code and note hashes.
The round-two record is
\path{../numerics/round-2/LOWER_CUTOFF_AND_PAIRED_EXTRACTION_CHECKS_20261017.json}.
The older remainder checker now labels its endpoint equality as an
equality; analytic strictness is supplied upstream.

Each normal run exits zero. The original failure injections remain
\texttt{--kernel-floor 9/40} and \texttt{--comparison-ceiling 13/20}.
The round-two injections are \texttt{--comparison-ceiling 69/100},
\texttt{--kernel-floor 3/10}, and \texttt{--derivative-allowance 3/2};
each exits one. The paired kernel floor is attained by a real zero at
horizontal distance $1/2$, including a double root counted twice.
For the disk estimate, with $r=1/2$,
\[
 F(w+r\zeta)=1+\frac7{10}\frac{\zeta-2/5}{1-(2/5)\zeta}
 \quad\Longrightarrow\quad |F'(w)/F(w)|=49/30>3/2.
\]
This is an exact disk-automorphism test, not a sample of $H_t$.
The comparison-ceiling injection fails a sufficient budget, not the
actual heat comparison. The separate class extremizers are not claimed
to be realized simultaneously by theta.

Finite arithmetic and class tests do not formally verify the classical
imports, contour shifts or infinite-domain analytic arguments. Their
purpose is authoring support. The current build record also distinguishes
Note 01's pre-existing single-link repair from its historical reviewed
hash; reversing that repair recovers the reviewed bytes. The notes and
reviews retain their source identities, and the earlier manuscript is
recoverable from commit \texttt{146b3f9}, without a new snapshot folder.

\section*{Acknowledgment of LLM assistance}
This draft was prepared for Edward Baker with substantial assistance
from OpenAI's \texttt{gpt-6.1-sol} model in Codex, using reasoning effort
\texttt{ultra}, on 17 October 2026. The assistance included derivation,
source comparison, same-model development work, numerical bookkeeping,
manuscript drafting and compilation checks. The research notes record
these contributions and their limitations. The baseline received Claude cross-model LLM review; cross-model review
of the round-two extension and specialist human review remain pending. Responsibility for checking,
revising and endorsing any eventual submission remains with its human
author.

\begin{thebibliography}{9}
\bibitem{polymath}
D. H. J. Polymath,
\emph{Effective approximation of heat flow evolution of the Riemann
$\xi$ function, and a new upper bound for the de Bruijn--Newman constant},
Research in the Mathematical Sciences \textbf{6}, article 31 (2019).
\href{https://doi.org/10.1007/s40687-019-0193-1}{doi:10.1007/s40687-019-0193-1}.
The source numbering used here follows
\href{https://arxiv.org/abs/1904.12438v2}{arXiv:1904.12438v2},
4 August 2019.
\bibitem{bellotti-wong}
C. Bellotti and P.-J. Wong,
\emph{Improved estimates for the argument and zero-counting function of
the Riemann zeta-function},
\href{https://arxiv.org/abs/2412.15470v2}{arXiv:2412.15470v2}, 7 July 2025,
Corollary 1.6. Used here only for a time-zero benchmark.
\end{thebibliography}
\end{document}
